\subsection*{MT-ART10. An integral Galois-module class forces a central zero}

Let $N/\mathbb Q$ be a tamely ramified $Q_8$-extension and let
$\chi$ be the unique irreducible complex character of degree two.
Fr\"ohlich's classical theorem identifies the two signs
\[
 W(\chi)=+1\quad\Longleftrightarrow\quad
       \mathcal O_N\text{ is free over }\mathbb Z[Q_8].
 \tag{MT-ART10a}
\]
The exact primary statement is Fr\"ohlich, \emph{The Galois
module structure of algebraic integer rings in fields with
generalised quaternion group} (1974), Theorem1, printed p.82,
with the locally free class invariant defined on pp.81--82.
The original proof is his \emph{Artin-root numbers and normal
integral bases for quaternion fields}, Inventiones17(1972),
143--166. Its full proof was not retrieved here; (MT-ART10a)
is used as the classical theorem, not claimed as a fresh proof.
Martinet's \emph{Fonctions z\^eta au point $1/2$}, \S VI,
printed114--115, gives the equivalent normal-basis formulation.

Consequently, if $\mathcal O_N$ is not free then
\[
 L(\tfrac12,\chi)=0,
 \qquad \operatorname{ord}_{s=1/2}\zeta_N(s)\ge2.
 \tag{MT-ART10b}
\]
These assertions are unconditional and are compatible with
both adopted GRH inputs.

\paragraph{Complete short deduction.}
Choose a cyclic subgroup $C_4\subset Q_8$ and a faithful
degree-one character $\lambda$ of $C_4$. Its induction is
$\chi$: its values are $2$ at the identity, $-2$ at the central
involution, and zero elsewhere, which is the irreducible
degree-two character. Artin induction-invariance therefore
identifies $L(s,\chi)$ with a nontrivial finite-order Hecke
function over $N^{C_4}$, so it is entire. Since $\chi$ is
selfdual, its completed function satisfies
\[
                 \Lambda(s,\chi)=W(\chi)\Lambda(1-s,\chi).
\]
Nonfreeness and (MT-ART10a) give $W(\chi)=-1$. Evaluation at
$s=1/2$ forces $\Lambda(1/2,\chi)=0$. The archimedean factors
and conductor power are finite and nonzero there, so the
finite-part $L$-function also vanishes. If $K=N^{\{1,-1\}}$
is the biquadratic subfield, the regular-character identity
gives the all-prime factorization
\[
                         \zeta_N(s)=\zeta_K(s)L(s,\chi)^2.
\]
The first factor is holomorphic at $1/2$, so the order is at
least two. More precisely the order of $L(s,\chi)$ there is
positive and odd, by the same functional equation. This proves
(MT-ART10b) without a general Artin-holomorphy hypothesis.

\paragraph{A completely specified classical example.}
Martinet, \emph{Modules sur l'alg\`ebre du groupe quaternionien}
(1971), TheoremIV.1 printed406 and Example3 printed407, gives
\[
 K=\mathbb Q(\sqrt5,\sqrt{41}),\qquad
 M=\frac{5+\sqrt5}{2}\frac{41+\sqrt{205}}{2},\qquad
 N=K(\sqrt M).
 \tag{MT-ART10c}
\]
This is a totally real tame $Q_8$-extension, ramified only at
$5$ and $41$, with nonfree integer ring. In his trace criterion
the two values are $205$ and $(1+5+41+205)/4=63$; they differ
modulo four, so the nonfree case applies. The radicand was
checked directly in the original page407 image, where the
second square root is $\sqrt{5\cdot41}$.
Thus this particular field has $\zeta_N(1/2)=0$ and satisfies
(MT-ART10b). For a convenient polynomial description, multiplying
$X^2-M$ over the four conjugates of $M$ gives
\[
 X^8-205X^6+13940X^4-378225X^2+3404025.
 \tag{MT-ART10d}
\]
An exact rational coefficient calculation of this norm is
recorded separately; the Galois group, ramification and integral
module classification are the cited classical example.

The programme bridge here is precise: tameness makes
$\mathcal O_N$ a locally free integral Galois module, and its
class in the locally free class group of $\mathbb Z[Q_8]$
controls one functional-equation sign. Rationalizing erases
this integral class while retaining the usual regular
representation. This is an example of the integral $K_0$
information distinguished in AG-LEFSCH-LATTICE and
AG-LEFSCH-DEFECT, not a general formula for all arithmetic
$K$-groups. It also shows that a residual condition such as
the nonvanishing of a Dedekind zeta function at $1/2$ cannot
be discarded merely because $M_H$ holds. It is not an example
inside the particular $S_5$ family of MT-S502; that family's
central nonvanishing gate remains its own arithmetic question.

